\documentclass[11pt]{article}
\usepackage[margin=1in]{geometry}
\usepackage{amsmath,amssymb,amsthm}
\usepackage{graphicx}
\usepackage[hidelinks]{hyperref}

\newtheorem{theorem}{Theorem}
\newtheorem{corollary}{Corollary}
\newtheorem{remark}{Remark}

\newcommand{\R}{\mathbb R}
\newcommand{\Sym}{\operatorname{Sym}}
\newcommand{\tr}{\operatorname{tr}}
\newcommand{\Div}{\operatorname{Div}}
\newcommand{\sgn}{\operatorname{sgn}}

\title{A smooth double-cone reduction for Question 4.12 of Guerra--Raita--Schrecker}
\author{Run \texttt{fa\_banach\_001}, lane 2}
\date{2026-07-03}

\begin{document}
\maketitle

\paragraph{Status.}
This is a partial result packet. It proves the smooth double-cone version of
Question 4.12 from Guerra--Raita--Schrecker, but it does not settle the rough
\(\mathrm{BV}^{\Div}\) formulation discussed in the same section.

\paragraph{Source.}
The source paper is Andre Guerra, Bogdan Raita, Matthew Schrecker,
\emph{Compensation phenomena for concentration effects via nonlinear elliptic
estimates}, arXiv:2112.10657. Question 4.12 appears on page 28 of the source
PDF:

\begin{center}
\includegraphics[width=0.95\linewidth]{figures/open_problem_crop.png}
\end{center}

\paragraph{Transcription of the target question.}
The question asks whether the inequalities of the two-dimensional section can
hold when the fields are constrained to the double cone
\[
  \Sym_2\cap\{\det\geq 0\}
  =\Sym_2^+\cup\Sym_2^- .
\]

\paragraph{Notation.}
If \(A=(a_{ij})\) is a \(2\times 2\) matrix field, let \(A_i\) be its \(i\)-th
row and let
\[
  d_i(A)=\operatorname{div} A_i=\partial_1 a_{i1}+\partial_2 a_{i2}.
\]
All matrix norms below are equivalent; in the proof we use the Frobenius norm.

\section*{Proof Intuition}

For a smooth field in the double cone, moving from the positive branch
\(\Sym_2^+\) to the negative branch \(\Sym_2^-\) can only happen through the
zero matrix, because trace zero inside the double cone forces both eigenvalues
to be zero. Therefore the branch-flip
\[
  B=\sgn(\tr A)A
\]
creates a positive semidefinite field without adding an interface divergence
term. Since \(\det(-A)=\det A\) in dimension two, the determinant is unchanged.
The source paper's positive-cone determinant inequality then applies to \(B\).

\section*{Main Result}

\begin{theorem}[Smooth double-cone determinant estimate]
Let \(A\in C_c^\infty(\R^2,\Sym_2^+\cup\Sym_2^-)\). Then
\[
  \int_{\R^2}\det A\,dx
  \leq
  \|d_1(A)\|_{L^1(\R^2)}\,\|d_2(A)\|_{L^1(\R^2)} .
\]
\end{theorem}

\begin{proof}
Put \(\tau=\tr A\), and define
\[
  B(x)=
  \begin{cases}
  A(x), & \tau(x)\geq 0,\\
  -A(x), & \tau(x)<0.
  \end{cases}
\]
Equivalently, \(B=\sgn(\tau)A\), with any value assigned to \(\sgn(0)\), since
\(\tau=0\) in \(\Sym_2^+\cup\Sym_2^-\) implies \(A=0\).
Thus \(B(x)\in\Sym_2^+\) for all \(x\), and \(\det B=\det A\).

We first record that \(B\) is Lipschitz. Indeed, the map
\(\Phi(M)=M\) on \(\Sym_2^+\) and \(\Phi(M)=-M\) on \(\Sym_2^-\) is Lipschitz
on the double cone. If two matrices lie on the same branch this is immediate.
If \(M,P\in\Sym_2^+\), then
\[
  \|M-P\|_F^2
  =\|M\|_F^2+\|P\|_F^2-2\operatorname{tr}(MP)
  \leq
  \|M\|_F^2+\|P\|_F^2+2\operatorname{tr}(MP)
  =\|M+P\|_F^2,
\]
because \(\operatorname{tr}(MP)\geq0\). This gives the cross-branch Lipschitz
bound.

It remains to identify the divergence. Let
\(\eta_\varepsilon(t)=t(t^2+\varepsilon^2)^{-1/2}\) and
\(B^\varepsilon=\eta_\varepsilon(\tau)A\). Then \(B^\varepsilon\to B\) in
\(L^1\), and
\[
  \partial_j B^\varepsilon_{ij}
  =
  \eta_\varepsilon(\tau)\,\partial_j A_{ij}
  +
  \eta_\varepsilon'(\tau) A_{ij}\,\partial_j\tau .
\]
For matrices in the double cone, \(|A_{ij}|\leq |\tau|\). Hence
\[
  |\eta_\varepsilon'(\tau) A_{ij}\partial_j\tau|
  \leq C|\partial_j\tau|,
\]
and the left-hand side tends pointwise to \(0\) wherever \(\tau\neq0\), while it
is zero where \(\tau=0\) after the factor \(|A_{ij}|\leq|\tau|\). Dominated
convergence gives
\[
  \partial_j B_{ij}=\sgn(\tau)\partial_j A_{ij}
\]
in distributions, and therefore
\[
  d_i(B)=\sgn(\tau)d_i(A),\qquad
  \|d_i(B)\|_1\leq \|d_i(A)\|_1 .
\]

The field \(B\) is compactly supported, Lipschitz, and positive semidefinite.
Since \(\Sym_2^+\subset\{a_{11},a_{22}\geq0\}\), Theorem 4.1 of the source
paper, with the rough-field extension in Remark 4.3, yields
\[
  \int_{\R^2}\det B\,dx
  \leq
  \|d_1(B)\|_1\,\|d_2(B)\|_1 .
\]
Using \(\det B=\det A\) and the divergence bound proves the claim.
\end{proof}

\begin{corollary}[Smooth \(L^2\) consequence]
Let \(\mathcal C\subset\Sym_2^+\cup\Sym_2^-\) be a closed cone such that
\[
  |M|^2\leq K\det M\qquad (M\in\mathcal C)
\]
for some \(K<\infty\). Then every
\(A\in C_c^\infty(\R^2,\mathcal C)\) satisfies
\[
  \|A\|_{L^2(\R^2)}
  \leq C_K \|\Div A\|_{L^1(\R^2)} .
\]
\end{corollary}

\begin{proof}
The cone assumption and the theorem give
\[
  \|A\|_2^2
  \leq K\int_{\R^2}\det A
  \leq K\|d_1(A)\|_1\|d_2(A)\|_1
  \leq C_K\|\Div A\|_1^2 .
\]
\end{proof}

\section*{Limitations and Review Notes}

The result should be read as a smooth answer. It does not resolve the rough
\(\mathrm{BV}^{\Div}\) problem discussed near Question 4.12. For a rough field
with a jump directly from a positive semidefinite value to a negative
semidefinite value, the transformation \(A\mapsto \sgn(\tr A)A\) creates an
additional jump term in the divergence. This is precisely the mechanism behind
the density obstruction in the source paper.

The main point for human review is the distributional identity
\[
  \Div(\sgn(\tr A)A)=\sgn(\tr A)\Div A
\]
for smooth double-cone-valued \(A\). The approximation by
\(\eta_\varepsilon(\tr A)\) is included above to make the vanishing of the
sign-derivative term explicit.

\section*{Novelty Check}

Bounded checks on 2026-07-03 searched the run indexes, the local parsed arXiv
corpus, and public web search for exact phrases including
\emph{Question 4.12}, \emph{double cone}, \(\Sym_2^+\cup\Sym_2^-\),
\emph{constrained Ornstein}, and the source arXiv id. No separate later answer
was found; exact hits led back to the source paper.

\begin{thebibliography}{9}
\bibitem{GRS}
A. Guerra, B. Raita, and M. Schrecker,
\emph{Compensation phenomena for concentration effects via nonlinear elliptic
estimates}, Ars Inveniendi Analytica (2024), Paper No. 1; arXiv:2112.10657.
\end{thebibliography}

\end{document}
